\section{气象匹配：省 $\times$ 年成分是不是初一特有的？}\label{app:weather}

正文第~\ref{sec:variance}~节把 30.1 个百分点的省 $\times$ 年成分同时归给“天气、区域传输、省级政策”，但当时没有气象数据，这项归因只是一个判断。本附录补上 ERA5 气象匹配，问三个问题：能复现多少？是不是初一特有？是不是纯天气？

\textbf{三个答案分别是：11.9 个百分点（39\%）、不是、不是。} 区域天气变量能复现省 $\times$ 年成分的约四成，但同样的解释力在除夕与前一天也在（11.8 与 9.7 个百分点），而且在省 $\times$ 年 \NOtwo{} 之上只剩 4.5 个百分点。\textbf{因此这个数不能读作“天气造成了初一尖峰”，它能支持的是一个更窄、但正文真正需要的结论：省 $\times$ 年成分是省 $\times$ 年尺度上的普遍背景，不是燃放日特有的，因而不能记在“税没对准初一”账上。}

\subsection{数据来源与构造}

气象取自 Open-Meteo 的 ERA5 存档接口（Hersbach et al., 2020），按城市经纬度点位插值，返回\textbf{本地日}（Asia/Shanghai）的日内统计量：10 m 风速日均值、2 m 气温日均值及日最高/最低值、地面气压日均值、相对湿度日均值。覆盖 306 城 $\times$ 12 个春节窗口 $\times$ 9 天（rel\_day $-7$ 至 $+1$），共 33{,}048 个城市-日；与 PM$_{2.5}$ 面板匹配后，进入估计的样本为 \textbf{3{,}663} 个城市-年（全样本 3{,}672 个，缺失 9 个，均因某一基线日缺测）。

\textbf{为什么只有两个窗口有边界层高度。} 边界层高度（BLH）是判断静稳天气最直接的变量。四条路都试过：Open-Meteo 的存档接口对 \texttt{boundary\_layer\_height} 返回全空值；historical-forecast 接口有该变量，但\textbf{只保留最近约两年}（2025、2026 两个窗口可得，2024 年 2 月为空）；NASA POWER 没有该参数；CDS 的 ERA5 需要注册与排队。因此本文用\textbf{10 m 风速}与\textbf{日温差}作通风条件的代理：风速越低、日温差越大，越可能静稳。\textbf{这是代理，不是 BLH 本身}；第~\ref{sec:blh-check}~节用两个可得窗口检验这个代理。

\textbf{一处已更正的坐标错配。} 初版城市坐标表把资阳、中山、大同分别错配到益阳、大连、大庆的坐标上。三处已改正并重抓对应城市-日。三个城市合计最多 36 个城市-年（占样本 1\%），\textbf{不足以改变任何量级}；正文全部数字都在改正后的数据上重算。

\subsection{检验设计}

沿用第~\ref{sec:variance}~节的序贯 $R^2$ 口径：先只用城市固定效应，再加入年份固定效应，最后加入省份 $\times$ 年份固定效应；相邻两步的 $R^2$ 增量就是该层吸收的方差份额。

\textbf{第一组变量：省级年均气象。} 该省该年所有可用城市的气象值的算术平均，8 个变量（风速、日温差、气压、湿度，各取当日值与基线六日均值）。它在构造上只随省份 $\times$ 年份变化，因此其 $R^2$ 增量\textbf{在数学上被省 $\times$ 年成分约束}，不会侵占城市 $\times$ 年份维度。

\textbf{第二组变量：城市相对省均的气象偏离。} 城市值减去所在省该年的均值，同样 8 个变量。它只随城市 $\times$ 年份变化，因此只能解释城市 $\times$ 年残差。两组相加恰好还原原始气象变量，所以“先省均、再偏离”与“直接放入 8 个气象变量”等价。

\textbf{两项反证。} 若 11.9 真的是“初一燃放遇到静稳天气”，它应当在初一显著大于相邻日。因此把同一套变量用在 rel\_day $-1$ 与 $0$ 上（同样以节前基线为参照）。另外，省级气象与省级工业活动同期相关，因此再加入省 $\times$ 年 \NOtwo{} 均值，看在活动之上气象还能解释多少。

\subsection{结果}

\begin{table}[H]
\centering
\caption{气象的增量解释力（$N=3{,}663$）}
\label{tab:weather-r2}
\small
\begin{tabular}{L{5.6cm}R{2.2cm}R{2.8cm}}
\toprule
\textbf{模型} & \textbf{累计 $R^2$} & \textbf{增量（百分点）} \\
\midrule
城市 + 年份固定效应 & 34.8\% & — \\
\quad + 省级年均气象（8 个变量） & 46.7\% & \textbf{11.9} \\
\quad + 城市相对省均的气象偏离 & 48.3\% & 1.6 \\
\midrule
城市 + 省份 $\times$ 年份固定效应 & 64.9\% & — \\
\bottomrule
\end{tabular}

\vspace{0.3em}
{\small 注：气象子样本自身的区域成分为 30.2 个百分点、城市 $\times$ 年残差为 35.1 个百分点，与正文全样本的 30.1 / 35.2 几乎相同。若改用调整 $R^2$（$K$ 个固定效应加回归元），区域成分为 28.4、气象增量为 12.9，份额升到 45\%；\textbf{39\% 对自由度口径敏感，不应读作精确值}。}
\end{table}

省级气象复现了区域成分的 39\%，本地气象偏离只解释城市 $\times$ 年残差的 5\%。但这两个数都不是本文想说的话，下面两项检验才是。

\begin{table}[H]
\centering
\caption{安慰剂“日”：同一套变量用在相邻日期上}
\label{tab:weather-day}
\small
\begin{tabular}{L{4.8cm}R{1.4cm}R{2.0cm}R{2.2cm}R{2.2cm}}
\toprule
\textbf{被解释变量} & \textbf{$N$} & \textbf{区域成分} & \textbf{气象增量} & \textbf{占区域成分} \\
\midrule
rel\_day $-1$（除夕前一天） & 3{,}663 & 43.8 & 9.7 & 22\% \\
rel\_day $0$（除夕） & 3{,}663 & 43.0 & 11.8 & 27\% \\
rel\_day $+1$（正月初一） & 3{,}663 & 30.2 & 11.9 & 39\% \\
\bottomrule
\end{tabular}
\end{table}

表~\ref{tab:weather-day}~是对“天气解释了初一尖峰”最直接的证伪。\textbf{三天的气象增量几乎一样（9.7—11.9 个百分点）}；初一的份额之所以更高，只是因为它的区域成分更小（30.2 对 43.0—43.8），不是分子更大。换句话说，这 11.9 度量的是\textbf{省 $\times$ 年尺度上气象与污染的普遍关联}，不是燃放日特有的天气放大。

\begin{table}[H]
\centering
\caption{在省 $\times$ 年 \NOtwo{} 之上，气象还能解释多少}
\label{tab:weather-no2}
\small
\begin{tabular}{L{7.4cm}R{2.6cm}}
\toprule
\textbf{模型（均在城市 + 年份 FE 之上）} & \textbf{增量（百分点）} \\
\midrule
只有省级气象（8 个变量） & 11.9 \\
只有省 $\times$ 年 \NOtwo{} 均值 & 13.0 \\
气象 + \NOtwo{} 的联合增量 & 17.5 \\
\midrule
\textbf{在 \NOtwo{} 之上，气象还能解释} & \textbf{4.5} \\
\bottomrule
\end{tabular}
\end{table}

表~\ref{tab:weather-no2}~说明第二件事：\textbf{气象增量里有约三分之二与省 $\times$ 年 \NOtwo{} 重叠}。\NOtwo{} 既是工业与交通活动的代理，也受气象影响，因此这一步\textbf{不能}分离“天气”与“区域活动”；它只说明 11.9 是“气象—活动”的联合可复现份额，\textbf{不是纯天气的贡献}。

\begin{table}[H]
\centering
\caption{省级年均气象的系数与聚类标准误（被解释变量：初一的 PM$_{2.5}$ 增量，\ugm）}
\label{tab:weather-coef}
\small
\begin{tabular}{L{5.0cm}R{2.4cm}R{2.4cm}R{1.8cm}}
\toprule
\textbf{变量} & \textbf{系数} & \textbf{聚类 SE} & \textbf{$t$} \\
\midrule
\multicolumn{4}{l}{\textit{初一当日}} \\
风速（m/s） & $-26.41$ & 4.79 & $-5.51$ \\
日温差（°C） & $+4.08$ & 0.52 & 7.82 \\
气压（hPa） & $-2.32$ & 0.56 & $-4.12$ \\
相对湿度（\%） & $+0.51$ & 0.28 & 1.79 \\
\midrule
\multicolumn{4}{l}{\textit{节前基线六日均值}} \\
风速（m/s） & $+22.11$ & 5.64 & 3.92 \\
日温差（°C） & $-5.11$ & 1.49 & $-3.44$ \\
气压（hPa） & $+1.61$ & 0.70 & 2.30 \\
相对湿度（\%） & $-1.28$ & 0.42 & $-3.02$ \\
\bottomrule
\end{tabular}

\vspace{0.3em}
{\small 注：标准误按省聚类（$G=31$），只做 $G/(G-1)$ 修正。\textbf{31 个集群下标准误本身不稳定，本表只用于判断量级，不作显著性结论。}}
\end{table}

表~\ref{tab:weather-coef}~还有一项规格检验。被解释变量是“当日减基线”的增量，所以同一天气变量在两天的系数理应\textbf{同号相反}。风速是 $-26.4$ 与 $+22.1$，日温差是 $+4.1$ 与 $-5.1$，气压是 $-2.3$ 与 $+1.6$——符号全部对得上，量级接近但不精确相等；相对湿度是 $+0.51$ 与 $-1.28$，符号对、量级不对称。这种模式符合“气象—污染关系”，但\textbf{不能证明因果}：任何在两天同向作用的省 $\times$ 年冲击也会产生对称系数。

\subsection{稳健性与读法}

\begin{table}[H]
\centering
\caption{稳健性}
\label{tab:weather-robust}
\small
\begin{tabular}{L{5.4cm}R{1.4cm}R{1.8cm}R{1.8cm}R{2.0cm}}
\toprule
\textbf{设定} & \textbf{$N$} & \textbf{区域成分} & \textbf{气象增量} & \textbf{占区域成分} \\
\midrule
基准（省级均值） & 3{,}663 & 30.2 & 11.9 & 39\% \\
留一均值（排除本市，$n\ge 2$） & 3{,}603 & 29.4 & 11.1 & 38\% \\
只用风速与日温差 & 3{,}663 & 30.2 & 7.9 & 26\% \\
\bottomrule
\end{tabular}
\end{table}

“省级均值”对单城市省份就是本市气象（北京、天津、上海、重庆、青海各只有 1 个面板城市）。表~\ref{tab:weather-robust}~的留一均值剔除本市重新计算，增量从 11.9 降到 11.1，结论不变。只用风速与日温差两个变量时降到 7.9，仍然可观。

\textbf{三条不能说的话。} 第一，不能说“区域天气解释了初一尖峰的 39\%”——正确说法是“8 个省级气象变量解释了 372 个省 $\times$ 年效应方差的约四成，且同样解释相邻日”。第二，不能把 11.9 与基线成分 21.0 相加或相减：11.9 是在增量上估的，同一套变量对基线水平只解释约 7 个点，两者重叠、不可加。第三，不能把“本地气象只解释 1.6”读成“剩下 95\% 不是气象”：缺少边界层高度，1.6 只是\textbf{可观测}本地气象的下界。

\textbf{能说的话只有一句，但足够。} 省 $\times$ 年成分在普通日子上同样存在，它由气象与区域活动共同构成，\textbf{不是燃放日特有的，更不是“税没对准初一”的证据}。这与正文的安慰剂结果（节前基线的区域成分仍有 21.0 个百分点）指向同一件事，只是路径不同。

\subsection{代理有效性：在两个有 BLH 的窗口上}\label{sec:blh-check}

BLH 只在 2025、2026 两个窗口可得，共 612 个城市-年。样本太小，不能重做正文的分解；但它足够回答一个更窄的问题：\textbf{风速与日温差是不是 BLH 的合格代理？}

\textbf{第一层，相关。} 省级夜间 BLH 与风速的相关系数为 0.60，与日温差为 $-0.42$；城市-日层面分别为 0.59 与 $-0.46$。方向全部符合物理直觉——风越大边界层越厚，晴夜温差越大边界层越薄——但相关并不紧。

\textbf{第二层，单独解释力。} 在城市 + 年份固定效应之上，8 个代理变量解释 8.2 个百分点；BLH 只用两个变量（当日与基线）就解释 5.0 个百分点。

\textbf{第三层，增量。} 这是关键。在 8 个代理之上再加 BLH，$R^2$ 只从 70.6\% 升到 70.7\%——\textbf{多 0.1 个百分点}。

\begin{table}[H]
\centering
\caption{代理与 BLH 的解释力（2025—2026，$N=612$）}
\label{tab:blh}
\small
\begin{tabular}{L{6.2cm}R{2.4cm}R{3.2cm}}
\toprule
\textbf{变量集} & \textbf{累计 $R^2$} & \textbf{增量（百分点）} \\
\midrule
城市 + 年份固定效应 & 62.4\% & — \\
\quad + 8 个代理变量 & 70.6\% & 8.2 \\
\quad + BLH（夜间，当日 + 基线） & 67.5\% & 5.0 \\
\quad + 8 个代理 + BLH & 70.7\% & 0.1（较 8 个代理） \\
\bottomrule
\end{tabular}

\vspace{0.3em}
{\small 注：只有 2 年、612 个观测。8 个回归元的自由度代价约 1.0 个百分点，所以“8 个代理”的增量含少量拟合噪声；BLH 只用 2 个回归元，与它的比较不受此影响。}
\end{table}

表~\ref{tab:blh}~的最后一行才是结论：BLH 单独解释 5 个点，\textbf{但在代理之上几乎不再增加}。代理不是 BLH，却携带了 BLH 能提供的几乎全部信息。因此正文用代理得到的 12 个点左右，不会因为“没有 BLH”而被系统性推翻；代理误差造成的衰减仍然存在，但它的幅度不能先验确定——这正是正文只写“可复现份额”、不写“天气贡献”的原因。

两处限制必须写明：相关与增量都只在 2 年、612 个观测上估计，不能外推到全部 12 个窗口；2025 与 2026 的天气若恰好比常年更平稳，代理的表现会好看一些。\textbf{这是一个校验，不是替换。}
